% ===================================================================
%  تابلوی شمارهٔ ۵ — خانه و ساختن آن.
%
%  Traduction, faite en 2026, en persan d'aujourd'hui, sur l'ido de
%  texto/io. Regles de la colonne : voir 10-tablo-01.tex. Une seule
%  numerotation, et le plus long tableau du livret : 69 blocs, 142
%  renvois.
%
%  CE TABLEAU EST UN DIALOGUE D'UN BOUT A L'AUTRE, ET LE PERSAN Y A
%  UNE DIFFICULTE QUE LE JAVANAIS AVAIT AUSSI, MAIS QUI NE SE REGLE
%  PAS DE LA MEME FACON. Le javanais a deux LEXIQUES paralleles —
%  ngoko et krama — et la colonne javanaise a du apprendre a
%  kolonoj.py que le niveau se choisit par qui parle a qui. Le persan
%  n'a pas deux lexiques : il a un jeu de verbes et de pronoms de
%  DEFERENCE — شما contre تو, فرمودن contre گفتن, تشریف بردن contre
%  رفتن, میل کردن contre خوردن —, et le partage se fait sur la
%  personne du verbe, non sur le vocabulaire des choses.
%
%  ET C'EST POUR CELA QUE kolonoj.py N'A RIEN A APPRENDRE ICI. Les dix
%  regles de la colonne persane lisent le CARACTERE — ی contre ي, ک
%  contre ك, les chiffres, le demi-espace —, et un caractere ne change
%  pas selon qui parle. Le controle du javanais devait distinguer le
%  recit du dialogue ; celui du persan traverse le tableau 5 sans que
%  la question se pose. Un controle qui lit l'octet est aveugle a la
%  politesse, et c'est ici une qualite.
%
%  LA COLONNE TIENT DONC L'ASYMETRIE SANS AIDE : Ioannes dit شما a son
%  oncle et a monsieur Blanko, il emploie می‌فرمایید et تشریف بیاورید ;
%  les adultes lui disent تو et emploient les formes ordinaires. C'est
%  ce qu'un livre persan pour enfants imprime, et l'alinea le demande
%  — l'oncle dit de Ioannes qu'il est serieux et qu'il faut tout lui
%  expliquer.
%
%  LE CHANTIER EST LE TABLEAU OU LE PERSAN EMPRUNTE LE MOINS DU
%  LIVRET, ET C'EST INATTENDU POUR UN VOCABULAIRE TECHNIQUE. آجر la
%  brique (61), گچ le platre (70), آهک la chaux (63), ماسه le sable
%  (62), سنگ la pierre, تیر la poutre (37), بام le toit (31), دیوار le
%  mur, پی les fondations (2), نردبان l'echelle, داربست l'echafaud
%  (110), بیل la pelle (82), کلنگ la pioche, تیشه l'herminette, اره la
%  scie (94), رنده le rabot (91), چکش le marteau (84), میخ le clou,
%  سندان l'enclume (127), انبر la tenaille (128), کوره la forge (126).
%  Vingt et un mots, tous iraniens, et pas un emprunt : on batit en
%  brique crue sur ce plateau depuis six mille ans, et la langue avait
%  tous les mots avant que quiconque arrive.
%
%  آجر EST LE PLUS ANCIEN DE TOUS, ET IL VIENT DE PLUS LOIN QUE LE
%  PERSAN : c'est l'akkadien agurru, la brique cuite, passe au vieux
%  persan puis a l'arabe (آجُرّ) puis retourne. Le mot a quatre mille
%  ans et il nomme la meme chose. Le releve n'a rien d'aussi vieux
%  ailleurs.
%
%  ET CE QUI EST EMPRUNTE DANS CE TABLEAU L'EST AU FRANCAIS, PAS A
%  L'ARABE : سیمان le ciment, بتن le beton, پلان le plan (76), آسانسور
%  l'ascenseur, پارکت le parquet (16), شومینه la cheminee. La couche
%  francaise du persan ne touche que ce qui est arrive au XIXe siecle
%  — la meme frontiere qu'aux tableaux 11 a 13 du javanais, ou le
%  neerlandais tient la ville moderne et le javanais tient la maison.
% ===================================================================

%%K t05-apar-1 apar
\VUsaut{6.0mm}
\VUtitre{15.0pt}{60}{تابلوی شمارهٔ ۵}
\VUsaut{1.4mm}
\VUtitre{11.4pt}{0}{خانه و ساختن آن.}
\VUsaut{2.2mm}

%%K t05-tit-1 sub
\VUpk{10.2pt}{40}{دیدار خوش.}
\VUsaut{3.0mm}

%%K t05-01-1 p
\textsc{یوهانس}. --- عمو جان، بسیار دوست دارم بدانم چگونه \VUgras{خانه}
می‌سازند.

%%K t05-01-2 p
\textsc{عمو} (80). --- این بسیار آسان است، دوست من ؛ با من بیا، خانه‌ای
را نشانت می‌دهم که آقای برناردوس (79) می‌سازد و من در آن خواهم نشست.

%%K t05-01-3 p
\textsc{یوهانس}. --- پس چه کسی \VUgras{پلان} \textsuperscript{(76)} این
خانه را کشیده است؟

%%K t05-01-4 p
\textsc{عمو}. --- \VUgras{معمار} \textsuperscript{(75)} ؛ او نقشهٔ بسیار
روشنی ارائه کرد که پذیرفته شد، و \VUgras{پیمانکار} \textsuperscript{(77)}
کارها را آغاز کرد.

%%K t05-01-5 p
\textsc{یوهانس}. --- پس پیمانکار کیست؟

%%K t05-01-6 p
\textsc{عمو}. --- او آقای بلانکو است، پدر دوست تو یاکوبوس. او با آقای
روفو، معمار، کارگران را سرپرستی می‌کند.

%%K t05-01-7 p
خوب! این هم آقای بلانکو که درست به هنگام با \VUgras{خط‌کش}
\textsuperscript{(78)} خود می‌آید، و آقای روفو با پلان خود.

%%K t05-01-8 p
روز به خیر، آقایان. حالتان چطور است؟

%%K t05-01-9 p
\textsc{آقای روفو}. --- بسیار خوب، سپاسگزارم. روز به خیر، یوهانس ؛ این
کودک چقدر بزرگ شده است!

%%K t05-01-10 p
\textsc{عمو}. --- آری، به راستی، او مرد کوچکی شده است. او بسیار جدی
است ؛ باید هر چه را می‌بیند برایش توضیح داد. امروز می‌خواهد بداند
چگونه خانه می‌سازند.

%%K t05-tit-2 sub
\VUpk{10.2pt}{40}{کارگران.}
\VUsaut{3.0mm}

%%K t05-01-11 p
\textsc{آقای بلانکو}. --- خوب، با من بیا، دوست کوچک ؛ همین حالا برایت
توضیح می‌دهم، زیرا کودکانی را که می‌خواهند بیاموزند بسیار دوست دارم.

%%K t05-01-12 p
آن کارگر را می‌بینی که \VUgras{بیل} \textsuperscript{(82)} به دست دارد :
او \VUgras{کارگر خاک‌برداری} \textsuperscript{(81)} است. او
\VUgras{کانال} \textsuperscript{(46)} می‌کند تا لولهٔ آب را بگذارند. او
زمین جایی را هم که خانه در آن است کنده است تا بنّا چهار \VUgras{دیوار}
را برپا کند. بخشی از این دیوارها که در زمین است \VUgras{پی}
\textsuperscript{(2)} نامیده می‌شود. فضایی که میان پی‌ها گرفته می‌شود
\VUgras{زیرزمین} \textsuperscript{(3)} را می‌سازد. این زیرزمین پنجرهٔ
کوچکی دارد که \VUgras{پنجرهٔ زیرزمین} \textsuperscript{(5)} نام دارد،
یک \VUgras{در} \textsuperscript{(4)} و یک پلکان.

%%K t05-01-13 p
\VUgras{بنّا} \textsuperscript{(108)} ساختن دیوارها را دنبال کرده و برای
\VUgras{درها} \textsuperscript{(8)} و \VUgras{پنجره‌ها}
\textsuperscript{(17)} جا گذاشته است.

%%K t05-01-14 p
\textsc{یوهانس}. --- اما آن بنّا را نمی‌بینم!

%%K t05-01-15 p
آ\textsuperscript{ی} \textsc{بلانکو}. --- اکنون کارش را در خانه به پایان
رسانده است. او نزدیک \VUgras{اسطبل} \textsuperscript{(54)}، بر
\VUgras{داربست} \textsuperscript{(110)} خود است ؛ او دیوار
\VUgras{رختشویخانه} \textsuperscript{(56)} را می‌سازد.

%%K t05-01-16 p
او ماله، تشتک و \VUgras{شاقول} \textsuperscript{(109)} دارد. اما این
نام‌ها را به یاد نخواهی سپرد.

%%K t05-01-17 p
\textsc{یوهانس}. --- البته که به یاد می‌سپارم، زیرا همین حالا از عمویم
مالهٔ کوچک و تشتکی خواهم خواست، و شاقول را خودم با ریسمانی خواهم ساخت.

%%K t05-tit-3 sub
\VUsaut{3.0mm}
\VUpk{10.2pt}{40}{مصالح.}
\VUsaut{3.0mm}

%%K t05-01-18 p
\textsc{عمو}. --- به‌به! بسیار خوب. اما به من بگو، یوهانس، برای ساختن
خانه به چه چیز نیاز است؟

%%K t05-01-19 p
\textsc{یوهانس}. --- به \VUgras{سنگ تراشیده} \textsuperscript{(59)} و
\VUgras{ملات} \textsuperscript{(65)} نیاز است.

%%K t05-01-20 p
\textsc{عمو}. --- درست است ؛ آن مرد را با کلاه بزرگ روی \VUgras{گاری}
\textsuperscript{(136)} خود ببین. او \VUgras{گاری‌ران}
\textsuperscript{(135)} است. هر روز \VUgras{قطعه‌های سنگ}
\textsuperscript{(60)} را به اینجا می‌آورد. او بار خود را در حیاط خالی
می‌کند، و \VUgras{سنگ‌تراشان} \textsuperscript{(83)} قطعه‌ها را
می‌تراشند و با \VUgras{ارهٔ سنگ} \textsuperscript{(85)} خود می‌برند.
\VUgras{کارگر ساده} \textsuperscript{(146)} آنها را نزد بنّا می‌برد که
جای‌شان می‌گذارد.

%%K t05-01-21 p
در همان هنگام، \VUgras{ملات‌ساز} \textsuperscript{(87)} ملات را آماده
می‌کند. او به \VUgras{ماسه} \textsuperscript{(62)}، \VUgras{آهک}
\textsuperscript{(63)} و \VUgras{آب} \textsuperscript{(64)} نیاز دارد.
آهک نزدیک او در \VUgras{کیسه} \textsuperscript{(71)} است ؛
\VUgras{شاگرد بنّا} \textsuperscript{(111)} ماسه را با \VUgras{فرغون}
\textsuperscript{(112)} برایش می‌آورد، و \VUgras{کارآموز}
\textsuperscript{(113)} کنار تلمبه (58) ایستاده است. او \VUgras{سطل}
\textsuperscript{(114)} را پر می‌کند. ملات به زودی آماده است.
\VUgras{ملات‌بر} \textsuperscript{(107)} آن را برمی‌دارد و با نردبان به
بالای داربست می‌برد، جایی که بنّایی این سوی خانه را به پایان می‌رساند.

%%K t05-01-22 p
و اکنون، کارگری که روی \VUgras{بام} \textsuperscript{(31)} کار می‌کند
چه نام دارد؟

%%K t05-01-23 p
\textsc{یوهانس}. --- او \VUgras{بام‌پوش} \textsuperscript{(118)} است.

%%K t05-tit-4 sub
\VUsaut{3.0mm}
\VUpk{10.2pt}{40}{بام خانه.}
\VUsaut{3.0mm}

%%K t05-01-24 p
آ\textsuperscript{ی} \textsc{بلانکو}. --- آری، اما نخست \VUgras{نجار}
\textsuperscript{(115)} می‌آید.

%%K t05-01-25 p
نجار \VUgras{خرپا} \textsuperscript{(35)} را می‌سازد. او \VUgras{تیرها}
\textsuperscript{(37)} را با \VUgras{میخ چوبی} \textsuperscript{(117)} به
هم می‌پیوندد. آن کارگران، آن بالا، با شلوارهای پهن مخملی خود، نجارند.
یکی تیر کوچکی را نگه داشته و دیگری میخ چوبی را با \VUgras{تبر}
\textsuperscript{(116)} خود می‌برد. آن که نزدیک \VUgras{دودکش}
\textsuperscript{(41)} است، بام‌پوش است. او چوب‌های باریک را روی (*)
تیرها میخ می‌کند، و \VUgras{سنگ لوح} \textsuperscript{(66)} را روی
چوب‌های باریک. گاهی سفال می‌گذارد.

%%K t05-noto-1 noto
\VUnotes{143.2mm}{%
(*) \VUgras{Balk-o} = آلمانی \textit{Balken}، انگلیسی \textit{joist}،
فرانسوی \textit{solive}، ایتالیایی \textit{travicello}، روسی
\textit{balka}، اسپانیایی و پرتغالی \textit{viga}.
ریشه‌ای فنی که تا کنون پذیرفته نشده است، اما شایستهٔ پیشنهاد است برای
برگرداندن معنای واژهٔ فرانسوی \textit{solive}، که نه trabo (فرانسوی
\textit{poutre}) است و نه trabeto. آن تیر چوبی چهارگوشی است که کف
تخته‌ای و سقف را نگه می‌دارد.%
}

%%K t05-01-26 p
بام اسطبل \VUgras{سفال‌پوش} \textsuperscript{(55)} است، بام خانه
\VUgras{لوح‌پوش} \textsuperscript{(32)} و بام ایوان از \VUgras{حلبی}
\textsuperscript{(24)}.

%%K t05-01-27 p
\textsc{یوهانس}. --- آیا بام‌پوش بام‌های حلبی را هم می‌سازد؟

%%K t05-01-28 p
آ\textsuperscript{ی} \textsc{بلانکو}. --- نه ؛ آن را \VUgras{حلبی‌ساز}
\textsuperscript{(121)} یا \VUgras{سرب‌کار} می‌سازد، همان که
\VUgras{پنجرهٔ بام} \textsuperscript{(38)} را روی بام خانه می‌گذارد. او
\VUgras{ناودان} \textsuperscript{(43)} و \VUgras{آبروی بام}
\textsuperscript{(42)} را هم می‌گذارد. او حلبی و ورق را لحیم می‌کند.
اوست که \VUgras{برق‌گیر} \textsuperscript{(40)} و \VUgras{بادنما}
\textsuperscript{(39)} را روی \VUgras{شیروانی} \textsuperscript{(36)}
نصب کرد.

%%K t05-01-29 p
\textsc{یوهانس}. --- پس خانه تمام شده است؟

%%K t05-tit-5 sub
\VUsaut{3.0mm}
\VUpk{10.2pt}{40}{ابزارها.}
\VUsaut{3.0mm}

%%K t05-01-30 p
آ\textsuperscript{ی} \textsc{بلانکو}. --- نه به تمامی. \VUgras{گچ‌کار}
\textsuperscript{(99)} با \VUgras{آجر} \textsuperscript{(61)}
دیوارهایی می‌سازد که خانه را به اتاق‌های گوناگون بخش می‌کنند. او با
\VUgras{گچ} \textsuperscript{(70)} \VUgras{سقف‌ها}
\textsuperscript{(26)} و دیوارها را می‌اندازد. اکنون سقف \VUgras{ایوان}
\textsuperscript{(33)} را کار می‌کند. او مانند بنّا \VUgras{ماله}
\textsuperscript{(100)} و \VUgras{تشتک} \textsuperscript{(101)} دارد.

%%K t05-01-31 p
سپس درودگر درها، پنجره‌ها، \VUgras{پارکت} \textsuperscript{(16)} و
\VUgras{کرکره‌ها} \textsuperscript{(19)} را می‌سازد.

%%K t05-01-32 p
اکنون در حیاط روی \VUgras{میز کار} \textsuperscript{(90)} خود کار
می‌کند. او \VUgras{اره} \textsuperscript{(94)} برای بریدن چوب (69)،
\VUgras{رنده} \textsuperscript{(91)}، \VUgras{گیره}
\textsuperscript{(92)}، \VUgras{اسکنه} \textsuperscript{(94 bis)}،
\VUgras{پرگار} \textsuperscript{(95)} و \VUgras{چکش} چوبی
\textsuperscript{(95 bis)} دارد.

%%K t05-01-33 p
\textsc{یوهانس}. --- به‌به! پس درودگری از بنّایی سرگرم‌کننده‌تر است.
عمو جان، برایم میز کار، اره و رنده خواهی خرید، زیرا به زودی برای مدور
لانه‌ای خواهم ساخت.

%%K t05-01-34 p
\textsc{عمو}. --- آری، پسرکم.

%%K t05-01-35 p
\textsc{یوهانس}. --- چه خرسندم! و آن که نزدیک درِ ورودی ایستاده است
کیست؟

%%K t05-tit-6 sub
\VUsaut{3.0mm}
\VUpk{10.2pt}{40}{رنگ زدن.}
\VUsaut{3.0mm}

%%K t05-01-36 p
آ\textsuperscript{ی} \textsc{بلانکو}. --- او \VUgras{نقاش ساختمان}
\textsuperscript{(102)} است. او روی نردبان خود با قوطی \VUgras{رنگ}
\textsuperscript{(103)} و \VUgras{قلم‌موی} \textsuperscript{(104)} خود
ایستاده است. او چوب درِ ورودی را رنگ می‌زند تا بیشتر بماند.

%%K t05-01-37 p
او کاغذ دیواری را هم در اتاق‌ها می‌چسباند.

%%K t05-01-38 p
\VUgras{پرده‌دوز} \textsuperscript{(107)} که روی \VUgras{نردبان}
\textsuperscript{(106)} بر \VUgras{بالکن} \textsuperscript{(28)} است،
\VUgras{کرکرهٔ پارچه‌ای} \textsuperscript{(29)} می‌گذارد.

%%K t05-01-39 p
کارگری که نزدیک پنجرهٔ بزرگ ایستاده است، \VUgras{شیشه‌بر}
\textsuperscript{(96)} است. او \VUgras{شیشه‌های پنجره}
\textsuperscript{(18)} را با \VUgras{بتونه} \textsuperscript{(73)} نصب
می‌کند. آن کارگر دیگر که نزدیک درِ زیرزمین زانو زده است،
\VUgras{قفل‌ساز} \textsuperscript{(97)} است. او \VUgras{قفل}
\textsuperscript{(6)} می‌گذارد. \VUgras{سوهانش} \textsuperscript{(98)}
نزدیک اوست.

%%K t05-01-40 p
\textsc{یوهانس}. --- آیا آن کارگری هم که با \VUgras{چکش}
\textsuperscript{(84)} خود بر آهن می‌کوبد، قفل‌ساز است؟

%%K t05-01-41 p
آ\textsuperscript{ی} \textsc{بلانکو}. --- درست‌تر بگویم، او آهنگر است
(125). او \VUgras{آهن‌آلات} \textsuperscript{(67)} را برای
\VUgras{برق‌کار} \textsuperscript{(124)} می‌سازد که روی بام
\VUgras{گاراژ} \textsuperscript{(51)} است. او \VUgras{کوره}
\textsuperscript{(126)}، \VUgras{سندان} \textsuperscript{(127)} و
\VUgras{انبر} \textsuperscript{(128)} دارد.

%%K t05-01-42 p
برق‌کار \VUgras{سیم} \VUgras{مسی} \textsuperscript{(44)} تلفن را نصب
می‌کند.

%%K t05-tit-7 sub
\VUpk{10.2pt}{40}{باغبانی.}
\VUsaut{3.0mm}

%%K t05-01-43 p
\textsc{عمو}. --- در انتهای \VUgras{حیاط} \textsuperscript{(45)}،
نزدیک \VUgras{نردهٔ آهنی} \textsuperscript{(47)} و \VUgras{دروازه}
\textsuperscript{(48)}، \VUgras{باغبان} \textsuperscript{(129)} را
می‌بینی. او \VUgras{باغچه} \textsuperscript{(49)} را آماده می‌کند. او
زمین را با \VUgras{بیل} \textsuperscript{(131)} خود بیل زده است، دانه
می‌پاشد و با \VUgras{شن‌کش} \textsuperscript{(130)} خود شن‌کش می‌کشد.
سپس با \VUgras{آب‌پاش} \textsuperscript{(132)} خود \VUgras{کرت‌ها}
\textsuperscript{(150)} را آب خواهد داد و گیاهان خواهند رویید.

%%K t05-01-44 p
\VUgras{مهتر} \textsuperscript{(133)} که همین حالا اسب را تیمار کرده و
به او خوراک داده است، \VUgras{کالسکه} \textsuperscript{(52)} را با
اسفنج، \VUgras{تلمبهٔ دستی} \textsuperscript{(134)} و سطلی آب پاک
می‌کند. سپس آن را دوباره به گاراژ خواهد برد و در را با \VUgras{کلون}
\textsuperscript{(53)} کلفت خواهد بست.

%%K t05-01-45 p
خوب، گمان می‌کنم خرسندی ؛ توضیح‌های بس داشتی.

%%K t05-01-46 p
\textsc{یوهانس}. --- آری، عمو جان، اما بسیار دوست دارم درون خانه را
ببینم.

%%K t05-01-47 p
\textsc{عمو}. --- این فردا خواهد بود. آقای روفو پلان را برایت توضیح
خواهد داد و نام هر اتاق را نشانت خواهد داد.

%%K t05-tit-8 sub
\VUsaut{3.0mm}
\VUpk{10.2pt}{40}{ترتیب درونی خانه.}
\VUsaut{2.4mm}

%%K t05-apar-2 apar
{\centering\textit{(پلان را ببینید.)}\par}
\VUsaut{2.4mm}

%%K t05-01-48 p
\textsc{معمار}. --- بیا، یوهانس، همین حالا تو را در بخش‌های گوناگون این
خانه می‌گردانم.

%%K t05-01-49 p
به \VUgras{طبقهٔ همکف} \textsuperscript{(7)} درآییم. این هم دکمهٔ
\VUgras{زنگ} \textsuperscript{(11)}. از سه \VUgras{پله}
\textsuperscript{(14)} \VUgras{ایوان ورودی} \textsuperscript{(9)} بالا
می‌رویم، از \VUgras{آستانهٔ} \textsuperscript{(10)} \VUgras{درِ ورودی}
\textsuperscript{(8)} می‌گذریم، و این هم که پای \VUgras{پلکان}
\textsuperscript{(13)} هستیم.

%%K t05-01-50 p
\textsc{یوهانس}. --- این راهروست.

%%K t05-01-51 p
\textsc{معمار}. --- نه، این \VUgras{سرسرا} \textsuperscript{(12)} است.
\VUgras{راهرو} \textsuperscript{(a)} در انتهاست. سمت راست، این هم
\VUgras{اتاق پذیرایی} \textsuperscript{(b)} و \VUgras{اتاق ناهارخوری}
\textsuperscript{(c)} ؛ روبه‌روی اتاق ناهارخوری \VUgras{آشپزخانه}
\textsuperscript{(d)} و \VUgras{اتاق خدمتکار} \textsuperscript{(e)}
هست، و سپس جلوتر \VUgras{اتاق کار} \textsuperscript{(f)}.
\VUgras{ایوان} \textsuperscript{(23)} با \VUgras{درِ شیشه‌ای}
\textsuperscript{(25)} خود نزدیک اتاق پذیرایی است.

%%K t05-01-52 p
اکنون از پلکان بالا خواهیم رفت. به \VUgras{نرده}
\textsuperscript{(15)} بگیر و پله‌ها را بشمار. چهارده تاست. این هم
\VUgras{راهروی بالا} \textsuperscript{(m)} ؛ ما در \VUgras{طبقهٔ}
\textsuperscript{(27)} اول هستیم.

%%K t05-tit-9 sub
\VUsaut{3.0mm}
\VUpk{10.2pt}{40}{طبقهٔ اول.}
\VUsaut{3.0mm}

%%K t05-01-53 p
روبه‌روی پلکان \VUgras{مستراح} \textsuperscript{(20)} و \VUgras{اتاق
آرایش} \textsuperscript{(j)} هست. سمت راست \VUgras{اتاق خواب}
\textsuperscript{(g)} زیبایی با دو پنجره رو به \VUgras{نمای}
\textsuperscript{(1)} خانه. سمت چپ \VUgras{حمام} \textsuperscript{(1)}
و \VUgras{اتاق مهمان} \textsuperscript{(i)} با \VUgras{شاه‌نشین}
\textsuperscript{(h)} بزرگ هست. نزدیک اتاق خواب \VUgras{اتاق کودکان}
\textsuperscript{(k)} است. آن رو به \VUgras{تراس} \textsuperscript{(21)}
پهناوری است. زیر این تراس مستراح دیگری (20) هست، و بر دیوار، این هم
\VUgras{زنگ} \textsuperscript{(22)} که برای شام به صدا درمی‌آید.

%%K t05-01-54 p
\textsc{یوهانس}. --- به \VUgras{طبقهٔ دوم} برویم.

%%K t05-01-55 p
\textsc{معمار}. --- طبقهٔ دومی نیست. زیر بام تنها \VUgras{اتاق‌های
زیرشیروانی} \textsuperscript{(33)} و انباری (34) هست. گردش ما به پایان
رسید ؛ می‌توانیم دوباره پایین برویم.

%%K t05-01-56 p
\textsc{یوهانس}. --- سپاسگزارم، آقا ؛ همهٔ اینها بسیار جالب بود. همین
حالا هر چه دیدم برای پدرم تعریف خواهم کرد.

\VUsaut{4.0mm}
\VUfilet{22mm}
